\documentclass[10pt]{article}

\usepackage[utf8]{inputenc}
\usepackage{CJKutf8}
\usepackage{geometry}
\geometry{a4paper, margin=0.75in}
\usepackage{booktabs}
\usepackage{enumitem}
\usepackage{setspace}
\usepackage{array}

\setlist{nosep, leftmargin=1.4em}
\renewcommand{\arraystretch}{1.15}

\begin{document}
\begin{CJK}{UTF8}{gbsn}

% ===================== COVER =====================
\begin{center}
    {\LARGE 北京师范大学}\\[4pt]
    {\large 硕士研究生学位论文开题报告}\\[2pt]
    {\large Beijing Normal University}\\[2pt]
    {\large Graduate Student Thesis Proposal Report}\\[18pt]
\end{center}

\begin{center}
\begin{tabular}{r p{0.5\linewidth}}
    姓名 Name: & Carlos Gustavo Salas Flores \\[4pt]
    学号 Student I.D.: & 202529030010 \\[4pt]
    导师 Supervisor: & Xu Hui \\[4pt]
    系所 School: & Business School \\[4pt]
    专业 Major: & World Economy and China \\[4pt]
    入学时间 Entry Date: & September 2025 \\[4pt]
\end{tabular}
\end{center}

\vspace{8pt}
\noindent\textbf{论文题目 Thesis Title:}\\
\emph{Economic Development through Political Materialism: How External Realignment, Not Regime Change, Drives Major Economic Disruption and Acceleration}

\newpage

% ===================== SECTION 1 =====================
\noindent{\large\bfseries 一、选题依据 Basis and Background of Topic Selection}

\vspace{6pt}
\noindent\textbf{Why this topic.} Development economics increasingly explains national prosperity through political variables, but those variables are usually measured by expert-coded indices (V-Dem, Polity V). Two findings motivate this thesis. Treier and Jackman (2008) show that Polity's democracy scores contain large measurement error, enough to reverse statistical conclusions. Little and Meng (2024) show that recent claims of global ``democratic backsliding'' are not supported by objective data (election results, competitiveness, journalists jailed or murdered): it is the experts' opinions that moved, not the countries. Any theory of development built on such indices inherits their errors.

\noindent\textbf{The gap in the dominant theory.} The most influential alternative, Acemoglu, Johnson and Robinson's (AJR) theory of \emph{inclusive} vs.\ \emph{extractive} institutions, explains why similar territories develop differently (their example: Nogales, Arizona vs.\ Nogales, Sonora). It has two weaknesses. First, \emph{circular measurement}: countries are classified as ``inclusive'' using the same outcomes (secure property rights, low expropriation) that the classification is supposed to explain. Second, a \emph{vague horizon}: when extractive economies grow (e.g., South Korea before 1987), the theory replies they must fail ``in the long run'' --- a claim with no time limit, which cannot be proven wrong.

\noindent\textbf{Research question.} When an economy breaks sharply from its previous path --- collapsing, like Venezuela (about $-75\%$ of GDP in 2013--2021) or the post-Soviet states ($-40\%$ to $-60\%$ in the 1990s), or taking off, like South Korea after 1961 or Poland after 1989 --- what actually causes the break? Is it domestic political change (what the literature loosely calls ``regime change''), or a change in the country's \emph{international position}: entering or losing a great power's sphere of influence?

\noindent\textbf{Hypotheses (in plain words).}
\begin{itemize}
    \item \textbf{H1:} Replacing rulers --- even changing the formal rules of the political game --- does not by itself produce a major economic break, as long as the effective organization of power and the country's international position stay the same.
    \item \textbf{H2:} Major economic breaks track changes in international position. Losing a patron (severance from a sphere of influence) brings \emph{disruption}; joining a richer sphere brings \emph{acceleration}.
    \item \textbf{H3:} Domestic political change coincides with major economic breaks only when the country's international position changes at the same time.
\end{itemize}

\noindent\textbf{Significance.} The framework explains cases the institutional theory cannot: North Korea's economic collapse in the 1990s with zero domestic political change (it lost the Soviet sphere), and South Korea's take-off from 1961 onward under a military regime (underwritten by its place in the US sphere). It also rebuilds the ``political variable'' in development economics without relying on contested expert-coded indices.

\newpage

% ===================== SECTION 2 =====================
\noindent{\large\bfseries 二、研究内容 Research Content}

\vspace{6pt}
\noindent\textbf{The framework in five questions.} The thesis describes any country through five questions: \emph{who governs}, \emph{how political power is organized}, \emph{through which permanent institutions power is exercised}, \emph{where the state stands internationally}, and \emph{what outcomes follow}. Key terms, in plain language:

\begin{center}\small
\begin{tabular}{p{0.28\linewidth} p{0.66\linewidth}}
\toprule
\textbf{Term} & \textbf{Plain meaning} \\
\midrule
Government & The people currently in office (president, ministers, ruling coalition). \\
Political regime & The formal rules for getting and keeping power (elections, succession). \\
Structure of political power & How power is \emph{actually} organized: whether powers are separated, who can control the rulers, party--state fusion (Garc\'ia-Trevijano). \\
State apparatus & The permanent machinery: army/police, bureaucracy, courts. \\
External alignment & The state's position toward great powers, alliances, and spheres of influence. \\
Sphere of influence & A region where an outside power has predominant influence, limiting members' freedom of action (Keal). \\
Major economic discontinuity & A sharp, sustained break in a country's economic path. Downward = \emph{disruption} (collapse, hyperinflation, capital flight). Upward = \emph{acceleration} (sustained take-off). \\
\bottomrule
\end{tabular}
\end{center}

A key distinction runs through the thesis: replacing rulers (or even changing formal rules) is not the same event as changing how power is actually organized; and a country changing its international position is not the same as the world power balance changing.

\noindent\textbf{What will be tested.} Hypotheses H1--H3 (Section 1), organized in a $3 \times 3$ comparative design: type of domestic political change (none / turnover / deep structural change) crossed with type of international change (none / severance from a sphere / entry into a richer sphere), with a signed outcome prediction for each cell.

\noindent\textbf{Case set.} Cases were surveyed and pre-classified (documented in \texttt{Cases.md}):

\begin{center}\small
\begin{tabular}{p{0.30\linewidth} p{0.30\linewidth} p{0.32\linewidth}}
\toprule
\textbf{Core (load-bearing)} & \textbf{Supporting} & \textbf{Hard cases (kept by design)} \\
\midrule
Cuba 1985--1995 (patron lost, zero domestic change) & South Korea 1961--1997 & Peru 1985--1992 (policy-driven crisis, no realignment) \\
Venezuela vs.\ Colombia 1998--2021 & Poland 1989--2004 (entry, acceleration) & Brexit 2016--2023 (partial exit, sub-threshold effects) \\
Post-Soviet cluster 1988--2004 (Belarus / Baltics / Ukraine / Central Asia) & Greece 2007--2015; Hungary 2010--2023 & Mexico 1994--1997 (entry with short crisis) \\
\bottomrule
\end{tabular}
\end{center}

\noindent\textbf{Thesis structure.} Ch.~1 Introduction; Ch.~2 Theoretical framework; Ch.~3 Literature review; Ch.~4 Research design and coding rules; Ch.~5--7 Empirical chapters (one per core case cluster); Ch.~8 Conclusion.

\newpage

% ===================== SECTION 3 =====================
\noindent{\large\bfseries 三、研究方法 Research Methods}

\vspace{6pt}
\noindent\textbf{Design.} Matched-pair comparison (most-similar systems): pairs of countries with similar material endowments (geography, resources, population) but different political/international histories --- e.g., Venezuela vs.\ Colombia, Belarus vs.\ the Baltics. Within each case, process tracing establishes the sequence: what changed first, what followed.

\noindent\textbf{Measuring the outcome.} \emph{Disruption} is coded with a growth-collapse criterion in the spirit of Rodrik (1999) (e.g., a deep, sustained output contraction); \emph{acceleration} with a symmetric growth-acceleration criterion in the spirit of Hausmann, Pritchett and Rodrik (2005). Both thresholds will be stated before coding the cases, and results will be re-run with alternative thresholds as a robustness check.

\noindent\textbf{Measuring the causes.} Domestic events are coded from the event taxonomy (turnover, coalition change, regime change, power-structure transformation, apparatus restructuring, policy change), using leader data (Archigos) and published coding rules. External alignment is coded from Keal's definition of a sphere of influence (predominant influence, constrained independence), using alliance data (ATOP/COW) and distinguishing \emph{severance} from \emph{entry}. Expert-coded democracy indices (V-Dem, Polity) are avoided as explanatory variables, given the measurement critique that motivates the thesis.

\noindent\textbf{Anticipated econometric concerns.}
\begin{itemize}
    \item \emph{Is N too small for causal claims?} The primary inference is comparative case analysis with explicit falsification criteria, not regression. As quantitative corroboration, synthetic control / interrupted time-series will be applied where donor pools exist (Venezuela, Cuba, Greece).
    \item \emph{Endogeneity / selection into alignment:} richer spheres tend to admit countries that are already succeeding (the Baltics reformed before joining the EU). Therefore the causal weight rests on \emph{severance} cases where the realignment was imposed from outside (the USSR disappeared; Cuba and North Korea did not choose it). Entry cases are used as confirmatory evidence, and a timing rule is enforced: outcomes must follow, not precede, the realignment they are attributed to.
    \item \emph{Confounders:} oil prices in Venezuela; the simultaneous systemic shock of 1989--91. Handled by matching (Colombia as oil-exporter control), within-cluster comparison (all post-Soviet states shared the systemic shock but differed in realignment), and by keeping policy-driven crises (Peru) as explicit scope conditions: H1 does not claim immunity from domestic policy shocks, only from changes of rulers and rules.
    \item \emph{Falsifiability:} the hypotheses fail if major breaks appear without realignment, if realignment produces no break, or if the sign is wrong (acceleration after severance, disruption after entry into a richer sphere).
\end{itemize}

\noindent\textbf{Data sources.} World Bank WDI; IMF WEO; Maddison Project (long-run GDP); EBRD Transition Reports; Archigos (leaders); ATOP/COW (alliances); national central banks (Venezuela, Cuba, with documented caveats).

\newpage

% ===================== SECTION 4 =====================
\noindent{\large\bfseries 四、研究目标 Research Objectives}

\vspace{6pt}
\noindent\textbf{Theoretical objective.} Rebuild the political variable in development economics: replace contested expert-coded indices with directly observable events (who rules, how power is organized, where the state stands internationally), and integrate the domestic structure of political power with the state's international position in a single framework.

\noindent\textbf{Empirical objective.} Test hypotheses H1--H3 on approximately 15 countries (1950--present), with published coding rules and a replication dataset, so that every classification (disruption/acceleration; severance/entry; turnover/structural change) can be checked and challenged.

\noindent\textbf{Expected contributions.}
\begin{enumerate}
    \item Explain cases the dominant institutional theory cannot: economic collapse without any domestic political change (Cuba, North Korea) and acceleration without prior democratization (South Korea after 1961).
    \item Reconnect development economics with international-relations realism: what happens to an economy depends on where the state sits in the international structure, and on what it gains or loses when that position changes.
    \item Clarify when political change matters economically: elections and even constitutional reforms may leave the economy untouched, while a change of patron or sphere can transform it --- with direct relevance for assessing small and middle powers aligned with or against great powers.
\end{enumerate}

\noindent\textbf{Work plan.}
\begin{center}\small
\begin{tabular}{p{0.22\linewidth} p{0.70\linewidth}}
\toprule
\textbf{Period} & \textbf{Task} \\
\midrule
Months 1--3 & Finalize theoretical framework and coding rules (with supervisor feedback). \\
Months 4--9 & Data collection and case coding; replication dataset. \\
Months 10--14 & Empirical chapters 5--7 (Cuba; Venezuela--Colombia; post-Soviet cluster). \\
Months 15--18 & Robustness checks, full draft, revision, defense preparation. \\
\bottomrule
\end{tabular}
\end{center}

\noindent\textbf{Feasibility.} All data sources are public; the case survey and pre-classification are already complete (\texttt{Cases.md}); the theoretical framework and hypotheses are written in full draft form. The main remaining work is coding and empirical analysis.

\vspace{24pt}

% ===================== SECTION 5 =====================
\noindent{\large\bfseries 五、导师意见 Tutor's Opinion}

\vspace{180pt}

\noindent 导师签名 Tutor Signature: \hrulefill\hspace{24pt}Date: \hrulefill

\end{CJK}
\end{document}
